\section{Welcome to Xeus-Haskell}


Xeus-Haskell is a Jupyter kernel for Haskell based on the native implementation of the Jupyter protocol, \href{https://github.com/jupyter-xeus/xeus}{xeus}.
\subsection{Preliminaries}


To start using or developing Xeus-Haskell, please note the following:
\begin{itemize}
\item \textbf{MicroHs Backend:} Xeus-Haskell uses \href{https://github.com/augustss/MicroHs}{MicroHs}, a minimal implementation of Haskell. Therefore, supported libraries and extensions are limited.
\item \textbf{Native Protocol:} Based on \href{https://github.com/jupyter-xeus/xeus}{xeus} (C++), so the kernel must be compiled for the specific platform you are using.
\item \textbf{WebAssembly Support:} You can compile this kernel to WebAssembly to run it directly in browsers!
\end{itemize}


\subsection{Quickstart}


Since this project is not yet available on Conda-forge or Emscripten-forge, you need to build it from source. We have organized the workflows using \href{https://pixi.sh}{pixi}.
\subsubsection{For Native JupyterLab}


\begin{minted}{text}
git clone https://github.com/jupyter-xeus/xeus-haskell
cd xeus-haskell
pixi run -e default prebuild
pixi run -e default build
pixi run -e default install
pixi run -e default serve # JupyterLab is ready!
\end{minted}


\subsubsection{For WebAssembly JupyterLite}


\begin{minted}{text}
git clone https://github.com/jupyter-xeus/xeus-haskell
cd xeus-haskell
pixi install -e prebuild
pixi run -e wasm-build prebuild
pixi run -e wasm-build build
pixi run -e wasm-build install
# pixi run -e wasm-build fix-emscripten-links # Run this if you encounter linking errors
pixi run -e wasm-build serve # JupyterLite is ready!
\end{minted}


\subsection{A Deep Dive into Xeus-Haskell}


\begin{itemize}
\item \hyperref[sec:contributing]{Contributing Guide}
\item \hyperref[sec:communication]{Communication Protocol}
\item \hyperref[sec:repl]{REPL Mechanism}
\item \hyperref[sec:todo]{Roadmap \& TODO}
\item \href{https://github.com/jupyter-xeus/xeus-haskell/blob/main/AGENTS.md}{AI Agent Guide}
\end{itemize}
